% (User input window)

\begin{remark}[Elliptic uniformization of the principal eigen-branch: Weierstrass \texorpdfstring{$\wp$}{wp} representation of \texorpdfstring{$\lambda_+(y)$}{lambda+(y)}]\label{rem:xi-terminal-zm-lambda-plus-wp-uniformization}
On the elliptic model of Conclusion \ref{con:xi-terminal-zm-elliptic-normalization}
\[
E:\ Y^2=X^3-X+\frac14
\]
the complex-analytic uniformization yields a period lattice \(\Lambda\subset\CC\) and a Weierstrass function \(\wp_\Lambda\) such that
\[
E(\CC)\cong \CC/\Lambda,\qquad
X=\wp_\Lambda(z),\qquad
Y=\frac{\wp_\Lambda'(z)}{2},
\]
and satisfying the differential equation
\[
\bigl(\wp_\Lambda'(z)\bigr)^2=4\wp_\Lambda(z)^3-4\wp_\Lambda(z)+1
\qquad(g_2=4,\ g_3=-1).
\]
Under this uniformization, the physical principal-branch weight function (Conclusion \ref{con:xi-terminal-zm-y-quadratic-extension-minpoly})
\[
y=X^2+Y-\frac12
\]
lifts to the elliptic function
\[
y(z)=\wp_\Lambda(z)^2+\frac{\wp_\Lambda'(z)}{2}-\frac12.
\]
Let \(P_1=(2,-\tfrac52)\in E(\QQ)\) (satisfying \(y(P_1)=1\), see Conclusion \ref{con:xi-terminal-zm-elliptic-rational-points-denominator-law}), and take \(z_1\in\CC/\Lambda\) such that
\(
P_1=(\wp_\Lambda(z_1),\wp_\Lambda'(z_1)/2)
\).
Because \(P_1\) is not a branch point of the cover \(\pi_y:E\to\mathbb P^1_y\) (Conclusion \ref{con:xi-terminal-zm-elliptic-y-hurwitz-passport}),
\(y(z)\) is locally invertible near \(z=z_1\); let \(z=z(y)\) denote a local inverse. Then the principal eigen-branch can be written as
\[
\lambda_+(y)=X\bigl(z(y)\bigr)=\wp_\Lambda\bigl(z(y)\bigr),
\qquad \lambda_+(1)=2.
\]
Hence \(\lambda_+(y)\) admits analytic continuation along any path avoiding the branch-value set, and a global representation is obtained by continuing \(z(y)\) on \(\CC/\Lambda\) and substituting back into \(\wp_\Lambda\).
Furthermore, from the differential identity in Conclusion \ref{con:xi-terminal-zm-leyang-ramification-critical-values}
\(
\dd y=(4XY+3X^2-1)\,\omega
\)
(where \(\omega=\dd X/(2Y)\)), we obtain: for any branch-cut-free region, take a local analytic section \(P(y)\in E(\CC)\) with \(\pi_y(P(y))=y\), and let
\(
X(y)=X(P(y))
\),\(
Y(y)=Y(P(y))
\), then
\[
\boxed{\;
\frac{\dd}{\dd y}\log \lambda_+(y)
=\frac{1}{X(y)}\frac{\dd X(y)}{\dd y}
=\frac{2Y(y)}{X(y)\,\bigl(4X(y)Y(y)+3X(y)^2-1\bigr)}\; }.
\]
Equivalently, on the minimal integral model \(E_0:\eta^2+\eta=x^3-x\), set \(x=\lambda\), \(\eta=y-\lambda^2\); then the same identity is
\[
\frac{\dd}{\dd y}\log \lambda(y)
=\frac{2\eta+1}{x\bigl(2x(2\eta+1)+3x^2-1\bigr)}.
\]
Therefore \(\bigl(\dd(\log\lambda_+)/\dd y\bigr)\dd y=\dd X/X\) is a third-kind Abelian differential on the elliptic curve, and its periods determine the multivalued analytic continuation of \(\log\lambda_+\).
Combined with the external anchoring in Remark \ref{rem:xi-terminal-elliptic-37a1-modular-anchor}, where \(E\) is modeled by \(X_0^+(37)\), these periods can also be interpreted in the period framework of level-\(37\) modular forms.
Its branch set and local types are fully determined by the ramification structure of \(\pi_y\): the finite branch values are
\(
y\in\{0,1\}\cup\{P_{\mathrm{LY}}(y)=0\}
\),
and each finite branch value has square-root-type Puiseux branching (Conclusion \ref{con:xi-terminal-zm-branch-puiseux-uniformization});
while at \(y=\infty\), \((y)_\infty=4O\) gives fourfold total ramification (Conclusion \ref{con:xi-terminal-zm-elliptic-y-hurwitz-passport}).
In particular, the nearest branch value on the real axis \(y_{\mathrm{LY}}\), and its associated degenerate eigenvalue \(\lambda_{\mathrm{LY}}\), are algebraically characterized by
\(
P_{\mathrm{LY}}(y_{\mathrm{LY}})=0
\)
and
\(
16\lambda_{\mathrm{LY}}^{3}-9\lambda_{\mathrm{LY}}^{2}+1=0
\)
respectively (Conclusion \ref{con:xi-terminal-zm-leyang-lambda-cubic}).
\end{remark}
